
% Expanded glossary entries for Houdini thesis
% Load with: \loadglsentries{glossary_full} before \begin{document}

% ===== Acronyms =====
\newacronym{vm}{VM}{Virtual Machine}
\newacronym{mac}{MAC}{Mandatory Access Control}
\newacronym{dac}{DAC}{Discretionary Access Control}
\newacronym{api}{API}{Application Programming Interface}
\newacronym{kvm}{KVM}{Kernel-based Virtual Machine}
\newacronym{qemu}{QEMU}{Quick emulator}
\newacronym{pid}{PID}{Process Identifier}

% ===== Core terms (define those used by others first) =====
\newglossaryentry{trick}{
  name=Trick,
  plural=Tricks,
  description={A scripted test case in Houdini that checks a specific aspect of container confinement, such as capabilities, policies, or escapes}
}
\newglossaryentry{runc}{
  name=Runc,
  description={The OCI-compliant low-level runtime used by Docker to spawn and run containers}
}
\newglossaryentry{ptrace}{
  name=Ptrace,
  description={A system call that allows one process to observe or control another; constrained by Yama and other security policies}
}
\newglossaryentry{lsm}{
  name=LSM,
  description={Linux Security Module; a kernel framework used by systems such as AppArmor, SELinux, and Yama}
}

% ===== Houdini terms =====
\newglossaryentry{Houdini}{
  name=Houdini,
  description={An execution framework for testing container confinement by running predefined \glspl{trick} in a controlled VM}
}
\newglossaryentry{overlay}{
  name=Overlay,
  description={A layered filesystem used in Houdini to inject Docker, Runc, and test files into the guest system}
}
\newglossaryentry{Flask}{
  name=Flask Server,
  description={The guest-side API in Houdini that receives requests from the host and orchestrates Trick execution}
}
\newglossaryentry{version pinning}{
  name=Version Pinning,
  description={Fixing exact versions of Docker, Runc, and the kernel to guarantee reproducible evaluation across time}
}
\newglossaryentry{regression testing}{
  name=Regression Testing,
  description={Re-running Tricks across versions to detect regressions in security enforcement}
}
\newglossaryentry{ablation testing}{
  name=Ablation Testing,
  description={Disabling one confinement mechanism at a time to measure its contribution to container security}
}

% ===== Container & runtime =====
\newglossaryentry{Docker}{
  name=Docker,
  description={A platform for building, running, and managing containers using an OCI runtime such as \gls{runc}}
}
\newglossaryentry{Docker Engine}{
  name=Docker Engine,
  description={The daemon that builds and runs containers; its version influences default confinement behaviors}
}
\newglossaryentry{container}{
  name=Container,
  description={A process running with isolated views of resources (via namespaces) and resource controls (via cgroups)}
}
\newglossaryentry{Dockerfile}{
  name=Dockerfile,
  description={A declarative build script that defines how to construct a container image}
}

% ===== Kernel isolation primitives =====
\newglossaryentry{namespaces}{
  name=Namespaces,
  description={Kernel features that isolate views of resources, including PID, network, mount, user, IPC, and UTS}
}
\newglossaryentry{AppArmor}{
  name=AppArmor,
  description={An \gls{lsm} that applies per-program security profiles to restrict file, capability, and network access}
}
\newglossaryentry{SELinux}{
  name=SELinux,
  description={An \gls{lsm} that enforces global label-based mandatory access control policies with fine-grained rules}
}
\newglossaryentry{Yama}{
  name=Yama,
  description={An \gls{lsm} that restricts use of \gls{ptrace} to harden process isolation}
}
\newglossaryentry{cgroups}{
  name=Cgroups,
  description={Linux control groups that limit and account for resources such as CPU, memory, and processes across containers}
}
\newglossaryentry{seccomp}{
  name=Seccomp,
  description={A kernel feature that filters and restricts system calls based on an allow- or deny-list}
}
\newglossaryentry{Linux capabilities}{
  name=Linux Capabilities,
  description={Fine-grained privileges split from root, such as \texttt{CAP\_NET\_BIND\_SERVICE} or \texttt{CAP\_SYS\_PTRACE}}
}
